% Appendix D, robustness of Section 6 (Oct 24, 2026). D.1: free entry, eta 1.5 (200-period arms, data_fi_torch/P15fi_*).
% D.2: slow stock (KMV permits), eta 0.63, 200-period fit (data_k200_torch/P063k_*). Numbers: output/model/experiments/
% housing_supply/20261021_paper_results/appD_values.json (p6_appendix.py); long run = stationary endpoint. At eta 1.5 the +10%
% supply run already matches credit on impact (1.585% vs 1.586%), so it is the matched experiment. Slow-stock matched supply (x1.4027,
% SUPMC) and credit for parents filled Oct 24 from data_k200_torch. Slow stock at eta 1.5 (P15k3_*, 300 periods) filled Oct 25 from
% appD_k15_values.json (appD_k15.py); matched = SUPM2 (+60%, births 2023 1.88% vs credit 1.92%).
% The previous version (old 100-period f3 runs) is in git (59d4c54e).
\providecommand{\pend}{\textcolor{red}{[pending]}}
\providecommand{\PRF}{../../output/model/experiments/housing_supply/20261021_paper_results}

\section{Robustness of the Policy Results}
\label{app:policy-robust}

This appendix repeats the experiments of Section~\ref{sec:policy} with a more
elastic supply and with a slowly adjusting stock.

\subsection{A more elastic supply}
\label{app:policy-eta}

Table~\ref{tab:app-policy-eta} repeats the experiments with free entry and
$\eta=1.5$. With a more elastic supply, a 10\% expansion lowers the price by
5.0\% on impact, against 8.6\% with $\eta=0.63$, and raises births in 2023 by
1.6\%, as much as the credit relaxation. It is therefore the matched
expansion. Credit again
moves births earlier: the completed fertility of the 2023 cohort falls by
1.0\%, while the supply expansion raises it by 2.6\%. In the steady state the price
is again the one at which fertility is at replacement, so credit for parents
raises it by about as much as with $\eta=0.63$, 34\%. With a more elastic
supply, more households fit at that price: 76\% more, against 36\%.

\begin{table}[H]
\centering
\caption{Counterfactuals from 2023 with $\eta=1.5$}
\label{tab:app-policy-eta}
\small
\begin{tabularx}{\textwidth}{@{}Xrrrcrr@{}}
\toprule
 & \multicolumn{3}{c}{Births} & Completed fertility & \multicolumn{2}{c}{House price} \\
\cmidrule(lr){2-4}\cmidrule(l){6-7}
 & 2023 & 2063 & Steady state & 2023 cohort & 2023 & Steady state \\
\midrule
Credit relaxation & $+1.6$ & $+0.1$ & $+4.8$ & $-1.0$ & $-0.1$ & $+2.6$ \\
Housing supply, $+10\%$ & $+1.6$ & $+4.3$ & $+10.0$ & $+2.6$ & $-5.0$ & $0.0$ \\
Credit for parents & $+5.1$ & $+14.2$ & $+75.6$ & $+5.8$ & $+0.1$ & $+33.9$ \\
\bottomrule
\end{tabularx}
\par\smallskip
\begin{minipage}{0.95\linewidth}
\footnotesize\textit{Notes:} As in Table~\ref{tab:policy-main}, with free
entry and $\eta=1.5$. The 10\% supply expansion raises births in 2023 as much
as the credit relaxation. Without policy, period fertility is 1.63 in 2023 and 1.67
in 2063.
\end{minipage}
\end{table}

\subsection{A slowly adjusting stock}
\label{app:policy-slow}

Here the housing stock adjusts slowly, as in
\citet{kaplanMitmanViolanteHousing2020}. Developers combine labor with building permits, which the government issues
in a fixed number each period. Construction and the stock then follow
\begin{equation}
 I_t=\delta_H H_0P_t^{\eta},\qquad H_t=(1-\delta_H)H_{t-1}+I_t .
 \label{eq:app-slow-stock}
\end{equation}
In a stationary equilibrium construction replaces depreciation, so the stock
lies on the same curve $H=H_0P^{\eta}$ as under free entry, and the 2007
equilibrium is the same. Along a transition only construction responds to the
price within a period, so the stock responds with an elasticity of about
$\delta_H\eta$, and the price absorbs most of a change in demand. The decline
in the preference for children from 2007 to 2023 is fitted under this rule, so
the 2023 state differs from the one in Section~\ref{sec:policy}. Everything
else is as in Section~\ref{sec:policy-design}.

Table~\ref{tab:app-policy-slow} reports the results with $\eta=0.63$. The credit
relaxation raises births by 1.7\% in 2023, but the gain falls to 0.6\% by 2063,
and the completed fertility of the 2023 cohort falls by 0.2\%. A 10\% supply
expansion lowers the price by 3.0\% in 2023 and raises births by only 0.4\%.
To raise births in 2023 as much as credit, the supply curve has to shift out
by 40\%. This matched expansion lowers the price by 10.9\% in 2023.
Construction keeps adding space, births are 12.2\% higher by 2063, and the
completed fertility of the 2023 cohort rises by 5.3\%. In the steady state the
price is back at its level without policy, so the supply expansions add as
many households as they add space, 10\% and 40\%, and credit adds 3.2\%.
Credit for parents raises births by 6.2\% in 2023 and by 32.9\% in the steady
state, when the price is 31.4\% higher.
Figure~\ref{fig:app-policy-slow-paths} shows the paths.

With $\eta=1.5$ the pattern is the same (lower panel of
Table~\ref{tab:app-policy-slow}). The credit relaxation raises births by 1.9\%
in 2023 and by 0.6\% in 2063, and the completed fertility of the 2023 cohort
falls by 0.8\%. To raise births in 2023 as much, the supply curve has to shift
out by about 60\%; this lowers the price by 10.8\% in 2023, raises births by
12.1\% by 2063 and the completed fertility of the 2023 cohort by 5.5\%.
Credit for parents raises births by 6.4\% in 2023 and the completed fertility
of the 2023 cohort by 7.5\%.

\begin{table}[H]
\centering
\caption{Counterfactuals from 2023 with a slowly adjusting stock}
\label{tab:app-policy-slow}
\small
\begin{tabularx}{\textwidth}{@{}Xrrrcrr@{}}
\toprule
 & \multicolumn{3}{c}{Births} & Completed fertility & \multicolumn{2}{c}{House price} \\
\cmidrule(lr){2-4}\cmidrule(l){6-7}
 & 2023 & 2063 & Steady state & 2023 cohort & 2023 & Steady state \\
\midrule
Credit relaxation & $+1.7$ & $+0.6$ & $+3.2$ & $-0.2$ & $+0.2$ & $+3.5$ \\
Housing supply, $+10\%$ & $+0.4$ & $+3.2$ & $+10.0$ & $+1.4$ & $-3.0$ & $0.0$ \\
Matched supply & $+1.7$ & $+12.2$ & $+40.3$ & $+5.3$ & $-10.9$ & $0.0$ \\
Credit for parents & $+6.2$ & $+17.1$ & $+32.9$ & $+5.9$ & $+1.8$ & $+31.4$ \\
\addlinespace
\multicolumn{7}{@{}l}{\emph{$\eta=1.5$}}\\
Credit relaxation & $+1.9$ & $+0.6$ & $+6.1$ & $-0.8$ & $-0.1$ & $+3.3$ \\
Housing supply, $+10\%$ & $+0.3$ & $+2.4$ & $+10.0$ & $+1.1$ & $-2.1$ & $0.0$ \\
Matched supply, $+60\%$ & $+1.9$ & $+12.1$ & $+60.0$ & $+5.5$ & $-10.8$ & $0.0$ \\
Credit for parents & $+6.4$ & $+17.2$ & $+64.7$ & $+7.5$ & $+1.1$ & $+29.8$ \\
\bottomrule
\end{tabularx}
\par\smallskip
\begin{minipage}{0.95\linewidth}
\footnotesize\textit{Notes:} As in Table~\ref{tab:policy-main}, $\eta=0.63$
(upper panel) and $\eta=1.5$ (lower panel), with the stock following
\eqref{eq:app-slow-stock}. Matched supply: the
supply shift that raises births in 2023 as much as the credit relaxation. Without
policy, period fertility is 1.64 in 2023 and 1.80 in 2063 with $\eta=0.63$, and
1.63 and 1.75 with $\eta=1.5$.
\end{minipage}
\end{table}

\begin{figure}[H]
\centering
\includegraphics[width=0.9\linewidth]{\PRF/fig_appD_paths_k200_eta063.pdf}
\caption{The credit relaxation and housing supply with a slowly adjusting stock}
\label{fig:app-policy-slow-paths}
\par\smallskip
\begin{minipage}{0.92\linewidth}
\footnotesize\textit{Notes:} \% differences from the no-policy path from the
same 2023 state, $\eta=0.63$, stock following \eqref{eq:app-slow-stock}.
Births in each four-year period (panel a), completed fertility of the cohort
aged 18--21 in the year shown (panel b), number of households (panel c) and
house price (panel d). Housing supply: the supply curve shifted out by 40\%,
which raises births in 2023 as much as the credit relaxation.
\end{minipage}
\end{figure}
